\chapter{Stimulant Use Disorder (Cocaine, Amphetamines)}
\index{Stimulant Use Disorder (Cocaine, Amphetamines)}
\label{sec:Stimulant Use Disorder (Cocaine, Amphetamines)}

\section{Overview}
Stimulant use disorder affects approximately 0.5--1\% of the population for cocaine and 0.2--0.4\% for amphetamine-type stimulants.
\section{Causes}
This condition happens because of cocaine blocking dopamine, norepinephrine, and serotonin transporters, and amphetamines additionally promoting neurotransmitter release from presynaptic vesicles, with chronic use leading to dopamine depletion, neurotoxicity, and extensive structural brain changes.     
\section{Risk Factors}

\begin{itemize}[noitemsep]
  \item Genetic predisposition and family history
  \item Advancing age
  \item Lifestyle factors (diet, exercise, smoking, alcohol)
  \item Environmental exposures
  \item Underlying medical conditions
  \item Medication use
\end{itemize}
\section{Signs and Symptoms}

\subsection{Common symptoms}
\begin{itemize}[noitemsep]
  \item \item You may experience a problematic pattern of stimulant use with criteria analogous to other substance use disorders. Withdrawal is marked by dysphoria, anhedonia, fatigue, hypersomnia, hyperphagia, and intense cravings. Pain or discomfort in the affected area    \item Pain or discomfort in the affected area
  \end{itemize}

\subsection{Emergency warning signs}
\begin{itemize}[noitemsep]
  \item Severe or worsening symptoms of stimulant use disorder (cocaine, amphetamines) require immediate medical evaluation
\end{itemize}
\section{Progression and Stages}
The condition may progress through different stages of severity over time. Early stages often involve mild or subtle symptoms that may be overlooked. Moderate stages involve more noticeable symptoms affecting daily function. Advanced stages may involve significant impairment and complications.
\section{Complications}
Medical complications include acute coronary syndrome, heart attack, involving bleeding stroke, rhabdomyolysis, seizures, high blood pressure, and for methamphetamine, severe dental decay (meth mouth), lung high blood pressure, and psychosis.

\begin{itemize}[noitemsep]
  \item Disease progression leading to worsening symptoms
  \item Secondary infections or injuries
  \item Side effects of treatment
  \item Reduced quality of life
  \item Emotional and psychological impact
\end{itemize}
\section{Diagnosis}
Doctors diagnose this using clinical interview. Comprehensive medical history and physical examination Laboratory tests to assess organ function and rule out other conditions Imaging studies to visualize affected structures Specialized testing depending on the affected system Consultation with relevant specialists Monitoring of disease progression over time

\begin{itemize}[noitemsep]
  \item Comprehensive medical history and physical examination
  \item Laboratory tests to assess organ function
  \item Imaging studies to visualize affected structures
  \item Specialized testing depending on the affected system
  \item Consultation with relevant specialists
\end{itemize}
\section{Prevention}
\begin{itemize}[noitemsep]
  \item Maintain a healthy lifestyle with balanced diet and exercise
  \item Avoid known risk factors
  \item Attend regular medical check-ups for early detection
  \item Follow recommended screening guidelines
\end{itemize}
\section{Treatment Options}
\begin{itemize}[noitemsep]
  \item Treatment options include behavioral interventions as first-line (CBT, contingency management), with no FDA-approved medications for stimulant use disorder. outlook: chronic relapsing course
  \item Methamphetamine use disorder has particularly high relapse rates.
  \item Medications to manage symptoms and address underlying mechanisms
  \item Lifestyle modifications and self-care measures
  \item Physical therapy and rehabilitation
  \item Surgical intervention when indicated
  \item Regular monitoring and follow-up care
\end{itemize}
\section{Living with It}
\begin{itemize}[noitemsep]
  \item Follow your treatment plan as prescribed
  \item Maintain regular follow-up with your healthcare provider
  \item Make necessary lifestyle adjustments
  \item Seek support from family, friends, or support groups
  \item Stay informed about your condition
\end{itemize}
\section{Outlook}
The outlook varies depending on the specific condition, severity at diagnosis, and response to treatment. Early intervention generally leads to better outcomes. Many conditions can be effectively managed with appropriate treatment. Regular follow-up care is important for monitoring and treatment adjustment.
